\section*{NeurIPS 论文核查表}

\begin{enumerate}
\item \textbf{论断}
\item[] 问题：摘要和引言中的主要论断是否准确反映论文贡献与范围？
\item[] 回答：\answerYes{}
\item[] 说明：摘要及第 1、7 节明确陈述了 RankMe 特有的宽度混淆、\SSRdag 的跨架构族范围、22 单元结果和主要局限。五任务两阶段结果被明确限定为概念验证，而非一般保证。

\item \textbf{局限}
\item[] 问题：论文是否讨论了本工作的局限？
\item[] 回答：\answerYes{}
\item[] 说明：第 7 节讨论了启发式参考选择、依赖模型池的性能、缺少形式化下游准确率保证、有限的同族证据和不完整的历史来源，两阶段流程只有五个任务的有限证据，以及保留 MV 候选间的预处理异质性。五任务诊断还包含遗憾更低的零标签流程。

\item \textbf{理论假设与证明}
\item[] 问题：每项理论结果是否给出完整假设与完整、正确的证明？
\item[] 回答：\answerYes{}
\item[] 说明：第 4 节给出 RankMe 结论所需的各向同性随机特征和长宽比假设，附录 A 使用 RankMe 的奇异值归一化给出 Marchenko--Pastur 推导。形式结论没有推广到任意旋转不变分数。

\item \textbf{实验结果可复现性}
\item[] 问题：论文是否充分披露复现主要实验结果及其结论所需的信息，无论是否提供代码和数据？
\item[] 回答：\answerNo{}
\item[] 说明：第 5 节给出有效秩、子空间对齐和 \SSRdag 的精确定义及算法。第 3 节和附录 B--G 记录支持集、VGG 提取设置、固定 LP 参数和缓存/脚本来源。保留结果表可核验汇总；部分原始缓存及历史环境未恢复，不能声称端到端历史复现。
\item \textbf{数据和代码的开放获取}
\item[] 问题：论文是否开放数据和代码，并提供足以忠实复现主要结果的说明？
\item[] 回答：\answerNo{}
\item[] 说明：配套仓库 \url{https://github.com/kenny67nju/crossarchitecturetransferability} 提供评分与分析代码、最终结果表、来源记录和核验说明。VGG 原始特征缓存不随仓库分发，部分历史缓存及环境不可用，因此尚未提供完整开放数据复现。

\item \textbf{实验设置与细节}
\item[] 问题：论文是否说明理解结果所需的全部训练和测试细节？
\item[] 回答：\answerYes{}
\item[] 说明：第 3 节描述模型池、线性探针协议、采样及逐候选预处理；第 5 节规定参考、$k$、裁剪、DiSCO 设置及汇总不确定性分析；第 6 节规定两阶段评估。附录 B--G 补充各池支持集、提取和 LP 设置、消融及有记录的计算环境。编码器保持冻结，仅拟合用于评估的线性探针。历史复现与计算资源记录的缺口见第 4、8 项。

\item \textbf{实验统计显著性}
\item[] 问题：论文是否恰当定义并报告误差条或其他统计显著性信息？
\item[] 回答：\answerYes{}
\item[] 说明：第 5 节的 22 单元主比较报告单元配对 bootstrap、按目标名称聚类的 $95\%$ 置信区间与双侧检验；聚类分析使用 $100{,}000$ 次重采样，并由配对 Wilcoxon 检验佐证。第 6 节报告五任务两阶段区间及其有限证据范围。

\item \textbf{实验计算资源}
\item[] 问题：论文是否为每项实验提供充分的计算资源、内存和运行时间信息？
\item[] 回答：\answerNo{}
\item[] 说明：附录~\ref{app:environment} 描述了NVIDIA L20 VGG 提取及有记录的华为昇腾 NPU 特征提取启动器、四设备分片配置及基于 CPU 的缓存特征分析。历史运行未一致保留精确硬件规格、软件版本快照和完整的逐次运行时间，因此现有信息仍不足以构成完整的计算资源报告。前向评估次数与实测运行时间加速比已作区分。

\item \textbf{伦理规范}
\item[] 问题：研究是否在各方面符合 NeurIPS 伦理规范 \url{https://neurips.cc/public/EthicsGuidelines}？
\item[] 回答：\answerYes{}
\item[] 说明：本工作在既有图像分类基准上评估预训练编码器并提出模型选择分数，不涉及人类受试者、私有数据或部署干预；作者已核查 NeurIPS 伦理规范。

\item \textbf{更广泛影响}
\item[] 问题：论文是否讨论工作的潜在正面与负面社会影响？
\item[] 回答：\answerNA{}
\item[] 说明：这是选择既有预训练编码器的基础方法，可能减少不必要的模型评估；除底层预训练模型和基准数据本身的影响外，它不引入直接、特定于应用的社会影响。

\item \textbf{安全保障}
\item[] 问题：对于具有高误用风险的数据或模型，论文是否说明了负责任发布所采用的保障措施？
\item[] 回答：\answerNA{}
\item[] 说明：发布内容是评分与分析代码及数值结果表；没有发布新的预训练模型、抓取数据集、生成系统或其他高风险资产。

\item \textbf{既有资产的许可}
\item[] 问题：论文是否适当标注所用代码、数据和模型的创作者，并明确说明和遵守许可及使用条款？
\item[] 回答：\answerNo{}
\item[] 说明：附录 B--E 引用了数据集原始论文或发布记录，包括派生数据集的底层来源。我们未提供涵盖全部数据集、模型权重和代码库的逐项许可证及使用条款清单，因此回答 No，而不声称许可文档完整。

\item \textbf{新资产}
\item[] 问题：论文引入的新资产是否有完善文档，并随资产一同提供？
\item[] 回答：\answerNo{}
\item[] 说明：评分代码、结果表、验证脚本和特征提取说明均在仓库中记录，但当前发布版本没有单独的软件许可证文件。

\item \textbf{众包与人类受试者研究}
\item[] 问题：若涉及众包或人类受试者，论文是否提供参与者说明全文、截图和报酬细节？
\item[] 回答：\answerNA{}
\item[] 说明：本研究未进行众包，也不涉及人类受试者。

\begin{samepage}
\item \textbf{机构审查委员会批准}
\item[] 问题：论文是否说明参与者风险、风险告知和 IRB 或所在国家/机构的等价审批？
\item[] 回答：\answerNA{}
\item[] 说明：本研究不涉及人类受试者或研究参与者，因此 IRB 批准与参与者风险披露不适用。
\end{samepage}

\begin{samepage}
\item \textbf{LLM 使用声明}
\item[] 问题：如果 LLM 是核心方法中重要、原创或非标准的组成部分，论文是否说明其使用？
\item[] 回答：\answerNA{}
\item[] 说明：LLM 不是所提出评分方法的组成部分。使用 AI 辅助进行文稿编辑、概念总览图制作及代码与审计支持；数学分数和报告的评估由所记录的数值程序计算，而非由 LLM 判断。作者对文本、图、代码和结果的核验承担责任。
\end{samepage}
\end{enumerate}
